\begin{apartats}

\apartat[1,25]{0,75}
Calcula el límit següent:
\[
\lim_{x\to-\infty}\frac{6x^2-5}{3x^2+2x}
\]

\begin{solucio}
Numerador i denominador tenen el mateix grau ($2$): el límit és el quocient dels
coeficients principals, $\dfrac{6}{3}=2$.
\end{solucio}

\begin{nomesllarg}
\apartat{0,75}
Calcula els límits següents:
\begin{graella}{2}
  \sa \lim_{x\to+\infty}\frac{1}{\sqrt[3]{x}} &
  \sa \lim_{x\to-\infty}\sqrt{1-x}
\end{graella}

\begin{solucio}
i) $\sqrt[3]{x}\to+\infty$, per tant el quocient tendeix a $0$.\\
ii) Si $x\to-\infty$, aleshores $1-x\to+\infty$ i $\sqrt{1-x}\to+\infty$.
\end{solucio}
\end{nomesllarg}

\apartat[1,25]{1}
Determina els nombres reals $a$ i $b$ perquè
\[
\lim_{x\to+\infty}\frac{ax^2+bx+1}{2x-1}=3 .
\]

\begin{solucio}
Si $a\neq0$, el numerador té grau $2$ i el denominador grau $1$: el límit seria infinit.
Per tant cal $\boxed{a=0}$.\\
Amb $a=0$, si $b\neq0$ numerador i denominador tenen grau $1$ i el límit val $\dfrac{b}{2}$
(si $b=0$ valdria $0$). Imposant $\dfrac b2=3$ s'obté $\boxed{b=6}$.
\end{solucio}

\end{apartats}
